%% Hand-authored circuitikz figure -- NOT generated by maestro2tex.
%% Topology transcribed device-by-device from the netlist of
%%   .simulation/castalia/tb_anatop_pixel_dacR2R12/maestro/results/maestro/
%%   Interactive.34/1/LoadRegTB/netlist/input.scs
%% (subckt anatop_pixel_dac_psu: 14 MOS + R0 + one INVX1).
%% Device dimensions are proprietary: only ratios and net names appear here.
\begin{figure}[htbp]
  \centering
  \resizebox{0.92\linewidth}{!}{%
\ctikzset{bipoles/length=0.9cm, transistors/scale=0.82, resistors/scale=0.85}
\begin{circuitikz}[
    every node/.append style={font=\footnotesize},
    gatelab/.style={font=\scriptsize, inner sep=1pt},
    devlab/.style={font=\scriptsize, text=black!60, inner sep=1pt},
    note/.style={font=\scriptsize, align=center, inner sep=2pt},
    railnode/.style={circle, fill=black, inner sep=1.05pt},
  ]

% ---- geometry -------------------------------------------------------
\def\xA{0} \def\xB{3.6} \def\xC{7.4} \def\xD{11.0}
\def\yvdd{10.2} \def\yvss{1.3}
\def\yP{7.4} \def\yrail{5.6} \def\yN{4.0} \def\yNb{2.2}

\draw[thick] (-1.7,\yvdd) -- (13.9,\yvdd) node[right]{$V_{\mathrm{DD}}$};
\draw[thick] (-1.7,\yvss) -- (13.9,\yvss) node[right]{$V_{\mathrm{SS}}$};

% ===== Branch A : PMOS diode over NMOS mirror -> SFBias ==============
\node[pmos](M3) at (\xA,\yP){};
\node[nmos](M5) at (\xA,\yN){};
\draw (M3.source) -- (\xA,\yvdd);
\draw (M3.drain) -- (\xA,\yrail);
\draw (\xA,\yrail) -- (M5.drain);
\draw (M5.source) -- (\xA,\yvss);
\node[railnode] at (\xA,\yrail){};
\node[gatelab, right=1pt, fill=white, inner sep=1.5pt] at (\xA,\yrail){\texttt{SFBias}};
\draw (M3.gate) -- ++(-0.35,0) |- (\xA,6.55);
\node[railnode] at (\xA,6.55){};
\node[gatelab, left] at (M5.gate){\texttt{A}};
\node[devlab, right=8pt] at (M3){M3};
\node[devlab, right=8pt] at (M5){M5};

% ===== Branch B : R0 + PMOS over NMOS diode -> A =====================
\draw (\xB,\yvdd) to[R, l_=$R_0$, *-] (\xB,8.35);
\node[railnode] at (\xB,8.35){};
\node[gatelab, right=3pt] at (\xB,8.35){\texttt{net6}};
\node[pmos](M2) at (\xB,\yP){};
\node[nmos](M4) at (\xB,\yN){};
\draw (M2.source) -- (\xB,8.35);
\draw (M2.drain) -- (\xB,\yrail);
\draw (\xB,\yrail) -- (M4.drain);
\draw (M4.source) -- (\xB,\yvss);
\node[railnode] at (\xB,\yrail){};
\node[gatelab, right=1pt, fill=white, inner sep=1.5pt] at (\xB,\yrail){\texttt{A}};
\node[gatelab, left] at (M2.gate){\texttt{SFBias}};
\draw (M4.gate) -- ++(-0.35,0) |- (\xB,4.85);
\node[railnode] at (\xB,4.85){};
\node[devlab, right=8pt] at (M2){M2};
\node[devlab, right=8pt] at (M4){M4};

% ===== Branch C : two stacked NMOS diodes -> B =======================
\node[pmos](M13) at (\xC,\yP){};
\node[nmos](M11) at (\xC,\yN){};
\node[nmos](M10) at (\xC,\yNb){};
\draw (M13.source) -- (\xC,\yvdd);
\draw (M13.drain) -- (\xC,\yrail);
\draw (\xC,\yrail) -- (M11.drain);
\draw (M11.source) -- (\xC,3.15);
\draw (\xC,3.15) -- (M10.drain);
\draw (M10.source) -- (\xC,\yvss);
\node[railnode] at (\xC,\yrail){};
\node[gatelab, right=1pt, fill=white, inner sep=1.5pt] at (\xC,\yrail){\texttt{B}};
\node[railnode] at (\xC,3.15){};
\node[gatelab, right=1pt, fill=white, inner sep=1.5pt] at (\xC,3.15){\texttt{C}};
\node[gatelab, left] at (M13.gate){\texttt{SFBias}};
\draw (M11.gate) -- ++(-0.35,0) |- (\xC,4.85);
\node[railnode] at (\xC,4.85){};
\draw (M10.gate) -- ++(-0.35,0) |- (\xC,3.15);
\node[devlab, right=8pt] at (M13){M13};
\node[devlab, right=8pt] at (M11){M11};
\node[devlab, right=8pt] at (M10){M10};

% ===== Branch D : current source + PMOS shunt -> VddDac ==============
\node[pmos](M0) at (\xD,\yP){};
\node[pmos](M1) at (\xD,\yN){};
\draw (M0.source) -- (\xD,\yvdd);
\draw (M0.drain) -- (\xD,\yrail);
\draw (\xD,\yrail) -- (M1.source);
\draw (M1.drain) -- (\xD,\yvss);
\node[railnode] at (\xD,\yrail){};
\node[gatelab, left] at (M0.gate){\texttt{SFBias}};
\node[gatelab, left] at (M1.gate){\texttt{B}};
\node[devlab, right=8pt] at (M0){M0};
\node[devlab, right=8pt] at (M1){M1};
\draw (\xD,\yrail) -- (13.0,\yrail);
\node[railnode] at (13.0,\yrail){};
\node[gatelab, above=3pt, align=center] at (13.0,\yrail){\texttt{VddDac}\\$\rightarrow$ bit drivers};

% ===== group labels =================================================
\node[note, anchor=north] at ({0.5*(\xA+\xB)},\yvss-0.4)
  {self-biased constant-$g_m$ reference:\\
   $R_0$ is 4 unit resistors in series and\\
   M2 is $4\times$ the width of M3, so $R_0$\\ alone sets the current};
\node[note, anchor=north] at ({0.5*(\xC+\xD)},\yvss-0.4)
  {shunt regulator:\\$V_{\texttt{VddDac}} = 2V_{GS,n} + |V_{GS,p}|$,\\
   M0 supplies, M1 sinks the surplus};

\end{circuitikz}
  }
  \caption[{Core of the DAC supply block.}]{Core of \texttt{anatop\_pixel\_dac\_psu}, the shunt regulator that generates the \texttt{VddDac} rail the ladder uses as its reference, transcribed from the \texttt{LoadRegTB} netlist. The self-biased constant-$g_m$ reference on the left sets \texttt{SFBias}; M13 drives that current through the stacked diodes M11 and M10 to hold \texttt{B} two gate--source voltages above $V_{\mathrm{SS}}$, and M1 sinks whatever of M0's current the ladder leaves, clamping \texttt{VddDac} one PMOS gate--source voltage above \texttt{B}: \SI{2.177}{\volt} unloaded from a \SI{2.5}{\volt} rail. There is no error amplifier, so load regulation is set by M1's $g_m$ alone. The five start-up and enable devices --- the weak pull-up and its two kick devices on \texttt{SU}, the shutdown pull-up on \texttt{SFBias} and the \texttt{En} inverter --- are omitted; none carries current once the loop is running. Device dimensions are proprietary and are not stated.}
  \label{fig:dacr2r12-psu}
\end{figure}
